\begin{table}[ht]
\centering
\renewcommand{\arraystretch}{1.4}
\renewcommand{\tabularxcolumn}[1]{m{#1}} 

\begin{tabularx}{\textwidth}{@{}ccc >{\centering\arraybackslash}X >{\centering\arraybackslash}m{1.8cm} @{}}
\toprule
\multicolumn{1}{c}{\textbf{Model}} & \multicolumn{1}{c}{\textbf{CAF}} & \multicolumn{1}{c}{\textbf{$p(n)$}} & \multicolumn{1}{c}{\textbf{$(N,p)_{N_g}$}} & \multicolumn{1}{c}{\textbf{Figures}} \\
\midrule

\multicolumn{5}{c}{\textit{without scalars}} \\
\midrule
\rowcolor{myDarkerBlue} GG & CAF & \eqref{eq:pNnoscalarGG} & $(5,0)_1$ &--- \\

\rowcolor{myLightBlue} BY & CAF & \eqref{eq:pNnoscalarBY} & $(3,0)_1$ &--- \\


\midrule
\multicolumn{5}{c}{\textit{with one scalar in the anti-fundamental representation - all on fixed flow}} \\
\midrule
\rowcolor{myDarkerBlue} & Not CAF: $N_g > 13$ & --- & --- & --- \\
\rowcolor{myDarkerBlue} \multirow{-2}{*}{GG} & CAF: $N_g\in\{1,\dots,13\}$ &\eqref{eq:pmaxfundNg} \& \eqref{eq:pminfundNgon} & $(5,6)_1$ $(5,5)_2$ $(5,4)_3$ $(5,3)_4$ $(5,2)_5$ $(5,0)_6$ $(5,0)_7$ $(5,0)_8$ $(5,0)_9$ $(5,0)_{10}$ $(5,0)_{11}$ $(5,0)_{12}$ $(5,0)_{13}$ &\ref{fig:cafgg_fund} \& \ref{fig:CAFGGNG9}  \\

\rowcolor{myLightBlue} & Not CAF: $N_g > 4$ & --- & --- & --- \\
\rowcolor{myLightBlue} \multirow{-2}{*}{BY} & CAF: $N_g\in\{1,\dots,4\}$  & \eqref{eq:pmaxfundNg}  & $(3,0)_1$ $(3,0)_2$ $(4,0)_3$ $(9,0)_4$ &\ref{fig:cafby_fund} \& \ref{fig:CAFNG9} \\


\midrule
\multicolumn{5}{c}{\textit{with one scalar in the anti-fundamental representation - Yukawa off fixed flow}} \\
\midrule
\rowcolor{myDarkerBlue} & Not CAF: $N_g > 13$ & --- & --- & --- \\
\rowcolor{myDarkerBlue} \multirow{-2}{*}{GG} &CAF: $N_g\in\{1,\dots,13\}$  & \eqref{eq:pmaxfundNg} \& \eqref{eq:pminfundNg}& $(5,22)_1$ $(5,20)_2$ $(5,18)_3$ $(5,16)_4$ $(5,14)_5$ $(5,12)_6$ $(5,10)_7$ $(5,8)_8$ $(5,6)_9$ $(5,4)_{10}$ $(5,2)_{11}$ $(5,0)_{12}$ $(5,0)_{13}$ &\ref{fig:cafgg_fund} \& \ref{fig:CAFGGNG9} \\

\rowcolor{myLightBlue} & Not CAF: $N_g > 4$ & --- & --- & --- \\
\rowcolor{myLightBlue} \multirow{-2}{*}{BY} & CAF: $N_g\in\{1,\dots,4\}$  &\eqref{eq:pmaxfundNg} \& \eqref{eq:pminfundNg} & $(3,10)_1$ $(3,4)_2$ $(4,0)_3$ $(9,0)_4$   &\ref{fig:cafby_fund} \& \ref{fig:CAFNG9} \\

\bottomrule
\end{tabularx}
\caption{Summary of CAF possibilities for the BY and GG models for different numbers of chiral families ($N_g$). The table shows the CAF status, the boundary functions $p(N,N_g)$ for vector-like fermions, and the minimal $N$ for each $N_g$ and the required minimal $p$ displayed as as $(N,p)_{N_g}$.\label{tab:conc_part1}}
\end{table}

Chiral gauge-Yukawa theories are the cornerstone of modern particle physics, exemplified by the Standard Model (SM). Even so, the understanding of their dynamics remains out of reach in various aspects. This thesis has closed some gaps in characterising their short-distance behaviour. We  considered the generalised Georgi--Glashow (GG) and the Bars--Yankielowicz (BY) models. Their high-energy dynamics was explored with a goal of realising complete asymptotic freedom (CAF) of the couplings. To evaluate physically relevant models, we added a single scalar in either the fundamental or the adjoint representation of $\SU(N)$. Exploiting the freedom to vary these parameters, we systematically evaluated the CAF conditions of the models, distinguishing between fixed-flow and off-fixed-flow regimes. The former describes the scenario where all the couplings vanishing at the same rate as the gauge coupling in the UV, while the latter describes the case where the Yukawa coupling vanishes faster than the gauge coupling. Through this analysis we varied the multiplicity of chiral families, as well as the vector-like fermions. The additional degrees of freedom originating from this later multiplicity proved necessary to realise CAF regimes in the models. 

Through this analysis, we discovered that the possible CAF regimes can be categorised only by the relation of the Yukawa and gauge couplings, and whether they exhibit "fixed-flow" or "off-fixed-flow" behaviour. The quartic coupling cannot be off-fixed-flow alone, since this will not be physically self-consistent, due to its beta functions term proportional to $\alpha_g^2$.

The CAF analysis demonstrates that all the considered classes of models can exhibit CAF regions within specific parameter spaces defined by the number of colours, chiral families, and vector-like families. These results are summarised in Table~\ref{tab:conc_part1} for models with no scalars or one fundamental scalar and in Table~\ref{tab:conc_part2} for models with an adjoint scalar. These tables detail the model type, the flow-trajectory type, the multiplicity of chiral families, the $p$-functions describing the upper and lower bounds on the needed vector-like fermions, the minimal models (in terms of the number of colours and required vector-like families for each viable chiral family), and lastly, references to related figures. 

This classification of the chiral gauge-Yukawa models provides a framework for identifying Grand Unified Theories (GUT)s that are UV complete. A model of particular interest is the $\SU(5)$ generalised Georgi--Glashow model with three chiral families. For this model, we found that when including one fundamental scalar, CAF regimes can be achieved by adding at least 4 vector-like fermions. This ensures fixed-flow CAF-dynamics, whereas 18 vector-like fermions are required for off-fixed-flow dynamics. In contrast, no such dynamics are found when the fundamental scalar is replaced with an adjoint scalar, in this case, the minimal model with three families would be $\SU(7)$ with at least 32 vector-like fermions. Recall, that a realistic GUT must include both a Higgs like scalar (represented here by the fundamental scalar) and a scalar that breaks the gauge group down to the Standard Model gauge sector, $\SU(3)_C\times\SU(2)_L\times \U(1)_Y$. Thus, the findings of this thesis serve as a preliminary step, as a complete, realistic model would require inclusion of both scalars. Such an analysis would be complicated by the extension of the Yukawa and quartic sector.

A complete understanding of chiral gauge-Yukawa models requires an analysis of their low-energy behaviour. While the conformal window has been investigated for purely fermionic models in \cite{Appelquist:2000qg,Pica:2016krb}, further analysis is needed to understand these dynamics when scalars are introduced. 


\begin{table}[ht]
\centering
\renewcommand{\arraystretch}{1.4}
\renewcommand{\tabularxcolumn}[1]{m{#1}} 

\begin{tabularx}{\textwidth}{@{}ccc >{\centering\arraybackslash}X >{\centering\arraybackslash}m{1.8cm} @{}}
\toprule
\multicolumn{1}{c}{\textbf{Model}} & \multicolumn{1}{c}{\textbf{CAF}} & \multicolumn{1}{c}{\textbf{$p(n)$}} & \multicolumn{1}{c}{\textbf{$(N,p)_{N_g}$}} & \multicolumn{1}{c}{\textbf{Figures}} \\
\midrule
\multicolumn{5}{c}{\textit{with one scalar in the adjoint representation - all on fixed flow}} \\
\midrule
\rowcolor{myDarkerBlue} & Not CAF: $N_g \in \mathbb{N} \setminus \{5,\dots,12\}$ &--- &--- &--- \\
\rowcolor{myDarkerBlue} \multirow{-2}{*}{GG} & CAF: $N_g\in\{5,\dots,12\}$ &\eqref{eq:padjscalar} \& \eqref{eq:Delta0} & $(13,2)_5$ $(9,1)_6$ $(7,1)_7$ $(6,1)_8$ $(6,1)_9$ $(6,1)_{10}$ $(5,1)_{11}$ $(5,1)_{12}$ & \ref{fig:final_combined_ng3}, \ref{fig:gg_combined_ng5} \& \ref{fig:gg_combined_ng6}\\

\rowcolor{myLightBlue} &Not CAF: $N_g \in \mathbb{N} \setminus \{3,4\}$ &--- &--- &--- \\
\rowcolor{myLightBlue} \multirow{-2}{*}{BY} & CAF: $N_g\in\{3,4\}$ & \eqref{eq:padjscalar} \& \eqref{eq:Delta0} &$(5,1)_3$ $(10,0)_4$ &\ref{fig:by_combined_ng5} \& \ref{fig:by_combined_ng6} \\


\midrule
\multicolumn{5}{c}{\textit{with one scalar in the adjoint representation - Yukawa coupling off fixed flow}} \\
\midrule
\rowcolor{myDarkerBlue} & Not CAF: $N_g > 9$ & --- & --- & --- \\
\rowcolor{myDarkerBlue} \multirow{-2}{*}{GG} & CAF: $N_g \in \{1,\dots,9\}$ & \eqref{eq:padjscalar} \& \eqref{eq:Delta0noY} & $(7,32)_1$ $(7,28)_2$ $(7,24)_3$ $(7,20)_4$ $(7,16)_5$ $(7,12)_6$ $(7,8)_7$ $(7,4)_8$ $(7,0)_9$ & \ref{fig:BYposroots}, \ref{fig:final_combined_ng2}, \ref{fig:final_combined_ng3}, \ref{fig:gg_combined_ng5} \& \ref{fig:gg_combined_ng6}\\

\rowcolor{myLightBlue} & Not CAF: $N_g > 4$ & --- & --- & --- \\
\rowcolor{myLightBlue} \multirow{-2}{*}{BY} & CAF: $N_g \in \{1,\dots,4\}$ & \eqref{eq:padjscalar} \& \eqref{eq:Delta0noY} & $(7,26)_1$ $(7,16)_2$ $(7,6)_3$ $(10,0)_4$ &  \ref{fig:GGadjNG37} \\


% \midrule
% \multicolumn{5}{c}{\textit{with one scalar in the adjoint representation - one quartic off fixed flow}} \\
% \midrule
% \rowcolor{myDarkerBlue} GG & Not CAF: $N_g\in\mathbb{N}$ & ---&--- & ---\\ 

% \rowcolor{myLightBlue} BY & Not CAF: $N_g\in\mathbb{N}$ &--- & ---&--- \\

\bottomrule
\end{tabularx}
\caption{Summary of CAF possibilities for the BY and GG models for different numbers of chiral families ($N_g$). The table shows the CAF status, the boundary functions $p(N,N_g)$ for vector-like fermions, and the minimal $N$ for each $N_g$ and the required minimal $p$ displayed as $(N,p)_{N_g}$. \label{tab:conc_part2}}
\end{table}
